\begin{table}[t]
\centering
\caption{Model-implied omitted-support bounds from monotone lower-tail completion}
\label{tab:omitted-support-bounds}
\footnotesize
\renewcommand{\arraystretch}{1.08}
\resizebox{0.98\textwidth}{!}{%
\begin{tabular}{>{\raggedright\arraybackslash}p{0.23\textwidth}cccccc}
\toprule
Target pair & Union share & $A$ share & $B$ share & Common-support $\Delta_{AB}$ & Monotone full-target interval & Min nonincreasing share \\
\midrule
Warm-up pair & 0.6411 & 0.2313 & 0.8414 & 0.00876 & [-0.10552,\ -0.07660] & 0.9933 \\
Nearby pair & 0.7472 & 0.3706 & 0.9338 & -0.00481 & [-0.02902,\ -0.02402] & 0.9978 \\
\bottomrule
\end{tabular}%
}
\begin{minipage}{0.98\textwidth}
\footnotesize Notes. The table tightens the generic common-support decomposition in Proposition~\ref{prop:common-support-selection} inside the maintained five-state hump calibration. For each PVEU target pair, the omitted cells all lie below the first supported liquid-asset point in the relevant age-state blocks under the verified current-cell map. We therefore complete the omitted mass using a monotone lower-tail restriction on the maintained point-summary normalized multiplier: within each relevant age-state block, $\eta$ is taken to be weakly decreasing in current liquid assets. The lower endpoint uses the first supported boundary value in each omitted block; the upper endpoint uses the maximum value over the omitted lower-tail in that block. The last column reports the minimum share of nonpositive asset-to-asset differences in the maintained point summary across the relevant age-state blocks. In this calibration those shares are at least $0.9933$ for the warm-up pair and at least $0.9978$ for the nearby pair, so the lower-tail monotonicity completion is numerically well aligned with the maintained point summary.
\end{minipage}
\end{table}
